\documentclass[a4paper,11pt]{article}

\usepackage[utf8]{inputenc}
\usepackage[T1]{fontenc}
\usepackage{amsmath,amssymb,amsthm}
\usepackage{physics}
\usepackage{hyperref}
\usepackage{geometry}
\geometry{margin=2.5cm}

\newcommand{\de}{\mathrm{d}}

\title{Funzione di Partizione a Temperatura Infinita \\[4pt]
       \large Vincolo sferico per spin complessi: $\sum_i |a_i|^2 = N$}
\author{}
\date{}

\begin{document}
\maketitle

%======================================================================
\section{Definizione del problema}
%======================================================================

Sia $a_i = \sigma_i + i\,\tau_i$ un numero complesso, con
$\sigma_i, \tau_i \in \mathbb{R}$, per $i = 1, \dots, N$.
Il modulo quadro è
\begin{equation}
  |a_i|^2 = \sigma_i^2 + \tau_i^2 .
\end{equation}
Vogliamo calcolare l'integrale
\begin{equation}\label{eq:Z}
  \mathcal{Z}
  \;=\; \int \prod_{i=1}^{N} \de\sigma_i\,\de\tau_i \;\;
        \delta\!\Bigl(N - \sum_{i=1}^{N} (\sigma_i^2 + \tau_i^2)\Bigr) .
\end{equation}
Questo è il volume dello spazio delle configurazioni a temperatura
infinita ($\beta = 0$) sotto il vincolo sferico
$\sum_i |a_i|^2 = N$.
Geometricamente, $\mathcal{Z}$ è l'area della superficie della sfera
$S^{2N-1}$ di raggio $R = \sqrt{N}$ in $\mathbb{R}^{2N}$.

%======================================================================
\section{Rappresentazione integrale della delta di Dirac}
%======================================================================

Il passo fondamentale, tipico della meccanica statistica, consiste
nell'usare la rappresentazione integrale della delta di Dirac.  Per ogni
$x \in \mathbb{R}$,
\begin{equation}\label{eq:delta_fourier}
  \delta(x) = \frac{1}{2\pi} \int_{-\infty}^{+\infty} \de\lambda \;
              e^{i\lambda x} .
\end{equation}
Equivalentemente, ponendo $z = c + i\lambda$ con $c > 0$ fissato
(contorno di Bromwich), si può scrivere
\begin{equation}\label{eq:delta_laplace}
  \delta(x) = \int_{c - i\infty}^{c + i\infty}
              \frac{\de z}{2\pi i}\; e^{z\,x} .
\end{equation}
Useremo la forma~\eqref{eq:delta_laplace} perché garantisce la
convergenza degli integrali gaussiani che seguiranno.

%======================================================================
\section{Inserimento della rappresentazione integrale}
%======================================================================

Sostituiamo l'Eq.~\eqref{eq:delta_laplace} nell'Eq.~\eqref{eq:Z},
con argomento $x = N - \sum_i(\sigma_i^2 + \tau_i^2)$:
\begin{equation}\label{eq:Z_step1}
  \mathcal{Z}
  = \int_{c - i\infty}^{c + i\infty} \frac{\de z}{2\pi i}\;
    \int \prod_{i=1}^{N} \de\sigma_i\,\de\tau_i \;\;
    \exp\!\Bigl[z\Bigl(N - \sum_{i=1}^{N}(\sigma_i^2 + \tau_i^2)\Bigr)\Bigr] .
\end{equation}
Separando nell'esponente la parte che dipende da $z$ e quella che
dipende dagli spin:
\begin{equation}\label{eq:Z_step2}
  \mathcal{Z}
  = \int_{c - i\infty}^{c + i\infty} \frac{\de z}{2\pi i}\;
    e^{zN}\;
    \prod_{i=1}^{N}
    \underbrace{\int_{-\infty}^{+\infty} \de\sigma_i
                \int_{-\infty}^{+\infty} \de\tau_i \;
                e^{-z(\sigma_i^2 + \tau_i^2)}}_{I_1(z)} .
\end{equation}
L'integrale si è \emph{fattorizzato} in $N$ contributi identici
$I_1(z)$: ciascuno è un integrale gaussiano bidimensionale nella
variabile complessa $a_i$.
Il contorno ha $\mathrm{Re}(z) = c > 0$, e ciò assicura che ogni
gaussiana convergaa.

%======================================================================
\section{Calcolo dell'integrale a singolo sito}
%======================================================================

L'integrale a singolo sito è
\begin{equation}
  I_1(z) = \int_{-\infty}^{+\infty} \de\sigma \int_{-\infty}^{+\infty}
           \de\tau \; e^{-z(\sigma^2 + \tau^2)} .
\end{equation}
Si fattorizza nelle due coordinate reali indipendenti:
\begin{equation}
  I_1(z) = \left(\int_{-\infty}^{+\infty} \de\sigma\; e^{-z\sigma^2}\right)
           \left(\int_{-\infty}^{+\infty} \de\tau\; e^{-z\tau^2}\right) .
\end{equation}
Ciascun fattore è un integrale gaussiano standard.
Per $\mathrm{Re}(z) > 0$:
\begin{equation}
  \int_{-\infty}^{+\infty} \de x \; e^{-zx^2}
  = \sqrt{\frac{\pi}{z}} .
\end{equation}
Pertanto:
\begin{equation}\label{eq:I1}
  I_1(z) = \sqrt{\frac{\pi}{z}} \cdot \sqrt{\frac{\pi}{z}}
         = \frac{\pi}{z} .
\end{equation}

%======================================================================
\section{Riduzione a un integrale su $z$}
%======================================================================

Sostituendo l'Eq.~\eqref{eq:I1} nell'Eq.~\eqref{eq:Z_step2}:
\begin{equation}
  \prod_{i=1}^{N} I_1(z) = \left(\frac{\pi}{z}\right)^{\!N} ,
\end{equation}
e dunque
\begin{equation}\label{eq:Z_contour}
  \mathcal{Z}
  = \pi^N \int_{c - i\infty}^{c + i\infty}
    \frac{\de z}{2\pi i}\; \frac{e^{Nz}}{z^N} .
\end{equation}
Rimane un singolo integrale nel piano complesso.
L'integrando $f(z) = e^{Nz}/z^N$ ha un \emph{polo di ordine $N$}
nell'origine $z = 0$, ed è analitico altrove.

%======================================================================
\section{Calcolo per residui}
%======================================================================

Il contorno d'integrazione è una retta verticale $\mathrm{Re}(z) = c > 0$
che racchiude il polo $z = 0$. Per il \textbf{teorema dei residui}:
\begin{equation}\label{eq:residue_thm}
  \int_{c - i\infty}^{c + i\infty} \frac{\de z}{2\pi i}\; \frac{e^{Nz}}{z^N}
  = \mathrm{Res}_{z=0}\; \frac{e^{Nz}}{z^N} .
\end{equation}

\paragraph{Sviluppo di Taylor.}
Espandiamo $e^{Nz}$ attorno a $z = 0$:
\begin{equation}
  e^{Nz} = \sum_{k=0}^{\infty} \frac{(Nz)^k}{k!}
          = \sum_{k=0}^{\infty} \frac{N^k}{k!}\, z^k .
\end{equation}

\paragraph{Divisione per $z^N$.}
\begin{equation}
  \frac{e^{Nz}}{z^N}
  = \sum_{k=0}^{\infty} \frac{N^k}{k!}\, z^{k - N} .
\end{equation}

\paragraph{Identificazione del residuo.}
Il residuo di un polo in $z = 0$ è il coefficiente di $z^{-1}$ nella
serie di Laurent. Il termine $z^{k-N}$ dà $z^{-1}$ quando
\begin{equation}
  k - N = -1 \qquad \Longrightarrow \qquad k = N - 1 .
\end{equation}
Dunque:
\begin{equation}\label{eq:residuo}
  \mathrm{Res}_{z=0}\; \frac{e^{Nz}}{z^N}
  = \frac{N^{N-1}}{(N-1)!} .
\end{equation}

\paragraph{Risultato.}
Sostituendo nell'Eq.~\eqref{eq:Z_contour}:
\begin{equation}\label{eq:Z_exact}
  \boxed{
  \mathcal{Z}
  = \frac{\pi^N \, N^{N-1}}{(N-1)!}
  = \frac{\pi^N \, N^{N-1}}{\Gamma(N)}
  }
\end{equation}
dove nell'ultimo passaggio abbiamo usato $\Gamma(N) = (N-1)!$ per
$N$ intero positivo.

%======================================================================
\section{Verifica: coordinate ipersferiche}
%======================================================================

Verifichiamo il risultato con un calcolo diretto.  Introduciamo le
coordinate ipersferiche in $\mathbb{R}^{2N}$, ponendo
\begin{equation}
  r^2 = \sum_{i=1}^{N} (\sigma_i^2 + \tau_i^2) ,
\end{equation}
cosicché l'elemento di volume si scrive
\begin{equation}
  \prod_{i=1}^{N} \de\sigma_i\,\de\tau_i
  = r^{2N-1}\,\de r\;\de\Omega_{2N-1} ,
\end{equation}
dove $\de\Omega_{2N-1}$ è l'elemento d'angolo solido sulla sfera
unitaria $S^{2N-1}$, il cui integrale totale è
\begin{equation}\label{eq:solid_angle}
  S_{2N-1} = \int \de\Omega_{2N-1}
           = \frac{2\,\pi^N}{\Gamma(N)} .
\end{equation}
L'integrale~\eqref{eq:Z} diventa:
\begin{equation}
  \mathcal{Z}
  = S_{2N-1} \int_0^{\infty} r^{2N-1}\,
    \delta(N - r^2)\,\de r .
\end{equation}

\paragraph{Calcolo della delta composita.}
Poiché $f(r) = N - r^2$ ha un unico zero $r_0 = \sqrt{N}$ per $r > 0$,
e $f'(r_0) = -2\sqrt{N}$, dalla regola per la delta di una funzione:
\begin{equation}
  \delta(N - r^2) = \frac{\delta(r - \sqrt{N})}{|{-2\sqrt{N}}|}
                  = \frac{\delta(r - \sqrt{N})}{2\sqrt{N}} .
\end{equation}
Sostituendo:
\begin{equation}
  \mathcal{Z}
  = S_{2N-1} \cdot
    \frac{(\sqrt{N})^{2N-1}}{2\sqrt{N}}
  = \frac{2\pi^N}{\Gamma(N)} \cdot
    \frac{N^{(2N-1)/2}}{2\,N^{1/2}} .
\end{equation}
Semplificando:
\begin{equation}
  \frac{N^{(2N-1)/2}}{N^{1/2}} = N^{(2N-2)/2} = N^{N-1} ,
\end{equation}
e dunque:
\begin{equation}
  \mathcal{Z}
  = \frac{\pi^N\, N^{N-1}}{\Gamma(N)} . \qquad \checkmark
\end{equation}
Il risultato coincide esattamente con l'Eq.~\eqref{eq:Z_exact}.

%======================================================================
\section{Entropia a temperatura infinita}
%======================================================================

L'entropia (microcanonica) a $\beta = 0$ è
\begin{equation}\label{eq:entropy}
  \boxed{
  S_\infty
  = \ln \mathcal{Z}
  = N \ln\pi + (N-1)\ln N - \ln\Gamma(N) .
  }
\end{equation}
Questa espressione è \emph{esatta} per ogni $N$ intero positivo.

%======================================================================
\section{Distribuzione congiunta delle intensità $I_i = |a_i|^2/N$}
%======================================================================

Definiamo per ciascun modo l'\emph{intensità normalizzata}
\begin{equation}\label{eq:Ii_def}
  I_i \;=\; \frac{|a_i|^2}{N} \;=\; \frac{\sigma_i^2 + \tau_i^2}{N} ,
  \qquad I_i \geq 0 .
\end{equation}
Il vincolo sferico $\sum_i |a_i|^2 = N$ diventa
\begin{equation}\label{eq:simplex}
  \sum_{i=1}^{N} I_i = 1 ,
\end{equation}
ossia le intensità vivono sul simplesso standard
$\Delta_{N-1} \subset \mathbb{R}^N$.

La distribuzione congiunta in $(\sigma_i, \tau_i)$ è
\begin{equation}\label{eq:joint}
  p(\{\sigma_i, \tau_i\})
  = \frac{1}{\mathcal{Z}}\;
    \delta\!\Bigl(N - \sum_{i=1}^{N} (\sigma_i^2 + \tau_i^2)\Bigr) .
\end{equation}
Vogliamo integrare sulle fasi e passare alle variabili $I_i$, ottenendo
$p(\vec{I}) = p(I_1, \dots, I_N)$.

%----------------------------------------------------------------------
\subsection{Cambio di variabili}
%----------------------------------------------------------------------

Per ciascun modo, passiamo a coordinate polari nel piano complesso:
\begin{equation}
  \sigma_i = \sqrt{N I_i}\,\cos\theta_i , \qquad
  \tau_i   = \sqrt{N I_i}\,\sin\theta_i ,
\end{equation}
con $I_i \geq 0$ e $\theta_i \in [0, 2\pi)$.  Lo Jacobiano è:
\begin{equation}
  \de\sigma_i\,\de\tau_i
  = r_i\,\de r_i\,\de\theta_i
  = \frac{1}{2}\,\de(r_i^2)\,\de\theta_i
  = \frac{N}{2}\,\de I_i\,\de\theta_i .
\end{equation}

%----------------------------------------------------------------------
\subsection{Integrazione sulle fasi}
%----------------------------------------------------------------------

Poiché la distribuzione~\eqref{eq:joint} non dipende dalle fasi
$\theta_i$, l'integrazione sulle $N$ fasi dà semplicemente un fattore
$(2\pi)^N$:
\begin{align}
  p(\vec{I})
  &= \int_0^{2\pi}\!\!\cdots\!\int_0^{2\pi}
     \prod_{i=1}^{N} \de\theta_i \;\;
     p(\{\sigma_i,\tau_i\}) \cdot
     \prod_{i=1}^{N} \frac{N}{2} \nonumber\\[6pt]
  &= \frac{(2\pi)^N}{\mathcal{Z}}
     \left(\frac{N}{2}\right)^{\!N}
     \delta\!\Bigl(N - N\sum_{i=1}^{N} I_i\Bigr) .
\end{align}
Usando $\delta(Nx) = \delta(x)/N$\,:
\begin{equation}\label{eq:pI_step}
  p(\vec{I})
  = \frac{(N\pi)^N}{N\,\mathcal{Z}}\;
    \delta\!\Bigl(1 - \sum_{i=1}^{N} I_i\Bigr) .
\end{equation}

%----------------------------------------------------------------------
\subsection{Sostituzione di $\mathcal{Z}$}
%----------------------------------------------------------------------

Dalla formula esatta~\eqref{eq:Z_exact},
$\mathcal{Z} = \pi^N N^{N-1}/(N-1)!$\,. Sostituendo:
\begin{equation}
  \frac{(N\pi)^N}{N\,\mathcal{Z}}
  = \frac{N^N\,\pi^N}{N} \cdot \frac{(N-1)!}{\pi^N\,N^{N-1}}
  = \frac{N^N \,(N-1)!}{N \cdot N^{N-1}}
  = (N-1)! .
\end{equation}
Pertanto:
\begin{equation}\label{eq:joint_I}
  \boxed{
  p(I_1, \dots, I_N)
  = (N-1)!\;\;\delta\!\Bigl(1 - \sum_{i=1}^{N} I_i\Bigr) ,
  \qquad I_i \geq 0 .
  }
\end{equation}
Le intensità sono \emph{uniformemente distribuite} sul simplesso
$\Delta_{N-1}$. Il fattore $(N-1)!$ è esattamente l'inverso del volume
del simplesso:
\begin{equation}
  \mathrm{Vol}(\Delta_{N-1})
  = \int_{\substack{I_i \geq 0}} \prod_{i=1}^{N} \de I_i \;\;
    \delta\!\Bigl(1 - \sum_{i} I_i\Bigr)
  = \frac{1}{(N-1)!} ,
\end{equation}
e dunque la normalizzazione è verificata.

Nella terminologia delle distribuzioni di probabilità,
l'Eq.~\eqref{eq:joint_I} è una distribuzione di \textbf{Dirichlet}
simmetrica con tutti i parametri uguali a $1$:
\begin{equation}
  \vec{I} \;\sim\; \mathrm{Dir}(\underbrace{1, 1, \dots, 1}_{N}) .
\end{equation}

%======================================================================
\section{Distribuzione marginale di $I_k$}
%======================================================================

Dalla distribuzione congiunta~\eqref{eq:joint_I}, la marginale di un
singolo modo $I_k$ si ottiene integrando su tutti gli altri:
\begin{equation}\label{eq:marginal_Ik_def}
  p(I_k)
  = (N-1)! \int_{\substack{I_j \geq 0 \\ j \neq k}}
    \prod_{j \neq k} \de I_j \;\;
    \delta\!\Bigl(1 - I_k - \sum_{j \neq k} I_j\Bigr) .
\end{equation}

%----------------------------------------------------------------------
\subsection{Volume del simplesso ridotto}
%----------------------------------------------------------------------

L'integrale residuo è sul simplesso in $N-1$ variabili
$\{I_j\}_{j \neq k}$ con il vincolo
$\sum_{j \neq k} I_j = 1 - I_k$.
Per $0 \leq I_k \leq 1$, questo è un simplesso $\Delta_{N-2}$ riscalato
di un fattore $(1 - I_k)$.
Il volume di un $(M{-}1)$-simplesso riscalato $\{x_j \geq 0,\,
\sum_j x_j = \Lambda\}$ in $M$ variabili è $\Lambda^{M-1}/(M-1)!$\,.
Qui $M = N - 1$ e $\Lambda = 1 - I_k$, dunque:
\begin{equation}\label{eq:simplex_volume}
  \int_{\substack{I_j \geq 0 \\ j \neq k}}
  \prod_{j \neq k} \de I_j \;\;
  \delta\!\Bigl(1 - I_k - \sum_{j\neq k} I_j\Bigr)
  = \frac{(1 - I_k)^{N-2}}{(N-2)!} .
\end{equation}

\paragraph{Dimostrazione.}
Lo si verifica per induzione, o direttamente con il cambio di variabile
$I_j = (1 - I_k)\,u_j$, che trasforma il vincolo in
$\sum_{j \neq k} u_j = 1$ con $u_j \geq 0$, e il Jacobiano è
$(1-I_k)^{N-2}$:
\begin{equation}
  \int \prod_{j \neq k} \de I_j \;\delta\!\Bigl(1 - I_k - \sum I_j\Bigr)
  = (1-I_k)^{N-2}
    \underbrace{\int \prod_{j \neq k} \de u_j \;\delta\!\Bigl(1 - \sum u_j\Bigr)}
    _{\mathrm{Vol}(\Delta_{N-2}) = 1/(N-2)!} .
\end{equation}

%----------------------------------------------------------------------
\subsection{Risultato}
%----------------------------------------------------------------------

Sostituendo l'Eq.~\eqref{eq:simplex_volume}
nell'Eq.~\eqref{eq:marginal_Ik_def}:
\begin{equation}
  p(I_k)
  = (N-1)! \cdot \frac{(1 - I_k)^{N-2}}{(N-2)!}
  = (N-1)\,(1 - I_k)^{N-2} .
\end{equation}
Dunque:
\begin{equation}\label{eq:marginal_Ik}
  \boxed{
  p(I_k) = (N-1)\,(1 - I_k)^{N-2} ,
  \qquad 0 \leq I_k \leq 1 .
  }
\end{equation}
Questa è una distribuzione $\mathrm{Beta}(1,\, N-1)$.

%----------------------------------------------------------------------
\subsection{Verifica della normalizzazione}
%----------------------------------------------------------------------

\begin{equation}
  \int_0^1 p(I_k)\,\de I_k
  = (N-1) \int_0^1 (1 - I_k)^{N-2}\,\de I_k
  = (N-1) \cdot \frac{1}{N-1} = 1 . \qquad \checkmark
\end{equation}

%----------------------------------------------------------------------
\subsection{Momenti}
%----------------------------------------------------------------------

\paragraph{Primo momento.}
\begin{align}
  \langle I_k \rangle
  &= (N-1) \int_0^1 I\,(1-I)^{N-2}\,\de I \nonumber\\[4pt]
  &= (N-1) \left[\frac{(1-I)^{N-1}}{-(N-1)}
     - \frac{(1-I)^{N}}{N}\right]_0^1
\end{align}
integrando per parti, oppure più direttamente con $u = 1-I$:
\begin{align}
  &= (N-1)\int_0^1 (1-u)\,u^{N-2}\,\de u
   = (N-1)\left[\frac{u^{N-1}}{N-1} - \frac{u^N}{N}\right]_0^1
   \nonumber\\[4pt]
  &= (N-1)\left(\frac{1}{N-1} - \frac{1}{N}\right)
   = 1 - \frac{N-1}{N} = \frac{1}{N} .
\end{align}
Dunque:
\begin{equation}
  \boxed{
  \langle I_k \rangle = \frac{1}{N} ,
  \qquad
  \langle |a_k|^2 \rangle = N\,\langle I_k \rangle = 1 ,
  }
\end{equation}
come deve essere per simmetria ($\sum_k I_k = 1$ e tutti i modi sono
equivalenti).

\paragraph{Secondo momento.}
Con $u = 1 - I$:
\begin{align}
  \langle I_k^2 \rangle
  &= (N-1)\int_0^1 (1-u)^2\,u^{N-2}\,\de u \nonumber\\[4pt]
  &= (N-1)\int_0^1 (1 - 2u + u^2)\,u^{N-2}\,\de u \nonumber\\[4pt]
  &= (N-1)\left[\frac{1}{N-1} - \frac{2}{N} + \frac{1}{N+1}\right]
   \nonumber\\[4pt]
  &= 1 - \frac{2(N-1)}{N} + \frac{N-1}{N+1}
   = \frac{N(N+1) - 2(N-1)(N+1) + N(N-1)}{N(N+1)} \nonumber\\[4pt]
  &= \frac{2}{N(N+1)} .
\end{align}

\paragraph{Varianza.}
\begin{equation}
  \boxed{
  \mathrm{Var}(I_k)
  = \langle I_k^2 \rangle - \langle I_k \rangle^2
  = \frac{2}{N(N+1)} - \frac{1}{N^2}
  = \frac{N-1}{N^2(N+1)} .
  }
\end{equation}

%----------------------------------------------------------------------
\subsection{Relazione con la distribuzione in $(\sigma_k, \tau_k)$}
%----------------------------------------------------------------------

Per completezza, riscriviamo la marginale anche in termini della
variabile originale $\rho = |a_k|^2 = N I_k$.
Poiché $\de\rho = N\,\de I_k$:
\begin{equation}
  p(\rho) = \frac{1}{N}\,p\!\left(I_k = \frac{\rho}{N}\right)
  = \frac{N-1}{N}\left(1 - \frac{\rho}{N}\right)^{\!N-2} ,
  \qquad 0 \leq \rho \leq N ,
\end{equation}
e in $(\sigma_k, \tau_k)$, integrando la fase
($\de\rho\,\de\theta/(2)$ con $\int_0^{2\pi}\de\theta = 2\pi$):
\begin{equation}
  p(\sigma_k, \tau_k)
  = \frac{N-1}{\pi N}\left(1 - \frac{\sigma_k^2+\tau_k^2}{N}\right)^{\!N-2} ,
  \qquad \sigma_k^2+\tau_k^2 \leq N .
\end{equation}

%----------------------------------------------------------------------
\subsection{Limite $N \to \infty$}
%----------------------------------------------------------------------

Per $N \to \infty$ a $\rho = N I_k$ fissato, $(1 - I_k)^{N-2}
= (1 - \rho/N)^{N-2} \to e^{-\rho}$, e le marginali diventano:
\begin{align}
  p(\sigma_k, \tau_k) &\;\xrightarrow{N\to\infty}\;
  \frac{1}{\pi}\,e^{-(\sigma_k^2 + \tau_k^2)} , \\[6pt]
  p(\rho) &\;\xrightarrow{N\to\infty}\; e^{-\rho} .
\end{align}
Ogni modo $a_k$ diventa una \emph{gaussiana complessa standard},
$a_k \sim \mathcal{CN}(0,1)$, e $I_k$ diventa un'esponenziale
riscalata: $N I_k \sim \mathrm{Exp}(1)$.

%======================================================================
\section{Interpretazione del moltiplicatore di Lagrange}
%======================================================================

Il parametro $z$ introdotto nella rappresentazione integrale della delta
gioca il ruolo di un \emph{moltiplicatore di Lagrange} (o potenziale
chimico) che impone il vincolo sferico.  Nella distribuzione a singolo
sito $\propto e^{-z(\sigma^2 + \tau^2)}$, il valore di aspettazione di
$|a_i|^2$ è:
\begin{equation}
  \langle |a_i|^2 \rangle_z
  = \frac{\displaystyle\int \de\sigma\,\de\tau\;
          (\sigma^2 + \tau^2)\,e^{-z(\sigma^2+\tau^2)}}
         {\displaystyle\int \de\sigma\,\de\tau\;
          e^{-z(\sigma^2+\tau^2)}}
  = -\frac{\partial}{\partial z} \ln I_1(z)
  = -\frac{\partial}{\partial z} \ln\frac{\pi}{z}
  = \frac{1}{z} .
\end{equation}
Richiedendo che $\sum_i \langle |a_i|^2 \rangle = N$, ossia
$N/z = N$, si ottiene $z = 1$. Questo è esattamente il valore del punto
sella dell'integrale~\eqref{eq:Z_contour}, ed è l'analogo del parametro
$\mu$ nel modello sferico di Berlin--Kac.

\end{document}
